% Source: lstm.tex
\begin{notechapter}
\chapter{Long Short-Term Memory}\label{ch:lstm}
\chaptersubtitle{The role of the cell state}

\section{Introduction}
Long Short-Term Memory (LSTM) networks are a type of recurrent
neural network (RNN) designed to address the \textbf{vanishing
gradient problem}, which plagues standard RNNs. The key innovation
in LSTMs is the introduction of the \textbf{cell state} (\(C_t\)),
which allows gradients to flow through time with minimal attenuation.
This document explains the LSTM architecture, focusing on the role
of the cell state in mitigating the vanishing gradient problem.

\section{The Vanishing Gradient Problem in RNNs}
In standard RNNs, the gradient of the loss with respect to the
weights is computed using \textbf{backpropagation through time (BPTT)}.
During this process, the gradient is multiplied by the same weight
matrix repeatedly as it propagates backward through time.

\subsection{Why Gradients Vanish}
The vanishing gradient problem arises because of the repeated
multiplication of the \textbf{hidden state weight matrix \(W_h\)}
during BPTT. If the eigenvalues of \(W_h\) are less than 1, the gradient
shrinks exponentially as it propagates backward. This makes it
difficult for the network to learn long-term dependencies. Conversely,
if the eigenvalues of \(W_h\) are greater than 1, the gradient can
explode, leading to the \textbf{exploding gradient problem}.

\subsection{Mathematical Explanation}
The hidden state \(h_t\) in an RNN is computed as:
\[
h_t = \sigma(W_h h_{t-1} + W_x x_t),
\]
where:
\begin{itemize}
    \item \(W_h\) is the weight matrix for the hidden state,
    \item \(W_x\) is the weight matrix for the input,
    \item \(\sigma\) is the activation function (e.g., tanh or sigmoid).
\end{itemize}

During BPTT, the gradient of the loss \(\mathcal{L}\) with respect
to the hidden state \(h_t\) is:
\[
\frac{\partial \mathcal{L}}{\partial h_t} = \frac{\partial \mathcal{L}}{\partial h_{t+1}} \cdot \frac{\partial h_{t+1}}{\partial h_t}.
\]
The term \(\frac{\partial h_{t+1}}{\partial h_t}\) depends on \(W_h\):
\[
\frac{\partial h_{t+1}}{\partial h_t} = \operatorname{diag}\big(\sigma'(W_h h_t + W_x x_{t+1})\big) \, W_h,
\]
where \(\sigma'\) is the derivative of the activation function. Over \(k\) steps the gradient is
multiplied by a product of \(k\) such Jacobians, whose norm is at most
\(\big(\gamma\,\|W_h\|_2\big)^k\), with \(\gamma = \sup|\sigma'|\) (\(1\) for \(\tanh\), \(1/4\) for the sigmoid)
and \(\|W_h\|_2\) the largest singular value. If \(\gamma\|W_h\|_2 < 1\), the gradient
\(\frac{\partial \mathcal{L}}{\partial h_t}\) is guaranteed to shrink exponentially
as it propagates backward through time. Conversely, a spectral radius of \(W_h\) greater than \(1\)
(largest eigenvalue modulus) is necessary for the gradient to explode (Pascanu et al., 2013). Note that it is the
singular values (or the spectral radius in the long run), together with \(\sigma'\), that matter: eigenvalues of
modulus smaller than 1 do not prevent transient growth for a non-normal \(W_h\).

\section{LSTM Architecture}
An LSTM consists of several components that work together to control
the flow of information. The main components are:
\begin{itemize}
    \item \textbf{Cell State} (\(C_t\)): The "memory" of the LSTM,
    which retains information over long time steps.
    \item \textbf{Hidden State} (\(h_t\)): The output of the LSTM
    at each time step.
    \item \textbf{Gates}: Mechanisms that control the flow of
    information into, out of, and within the cell state.
\end{itemize}

\subsection{Equations of an LSTM}
The LSTM updates its cell state and hidden state at each time step
using the following equations:

\subsubsection{Forget Gate (\(f_t\))}
The forget gate determines what information to discard from the
previous cell state (\(C_{t-1}\)):
\[
f_t = \sigma(W_f \cdot [h_{t-1}, x_t] + b_f),
\]
where:
\begin{itemize}
    \item \(W_f\) is the weight matrix for the forget gate,
    \item \(h_{t-1}\) is the hidden state at the previous time step,
    \item \(x_t\) is the input at the current time step,
    \item \(b_f\) is the bias term for the forget gate,
    \item \(\sigma\) is the sigmoid activation function.
\end{itemize}

\subsubsection{Input Gate (\(i_t\)) and Candidate Cell State (\(\tilde{C}_t\))}
The input gate determines what new information to add to the cell
state, and the candidate cell state (\(\tilde{C}_t\)) represents
the new information:
\[
i_t = \sigma(W_i \cdot [h_{t-1}, x_t] + b_i),
\]
\[
\tilde{C}_t = \tanh(W_C \cdot [h_{t-1}, x_t] + b_C),
\]
where:
\begin{itemize}
    \item \(W_i\) and \(W_C\) are the weight matrices for the input
    gate and candidate cell state,
    \item \(b_i\) and \(b_C\) are the bias terms,
    \item \(\tanh\) is the hyperbolic tangent activation function.
\end{itemize}

\subsubsection{Cell State Update (\(C_t\))}
The cell state is updated by combining the previous cell state
(\(C_{t-1}\)) and the candidate cell state (\(\tilde{C}_t\)),
controlled by the forget gate (\(f_t\)) and input gate (\(i_t\)):
\[
C_t = f_t \odot C_{t-1} + i_t \odot \tilde{C}_t,
\]
where \(\odot\) denotes element-wise multiplication.

\subsubsection{Output Gate (\(o_t\)) and Hidden State (\(h_t\))}
The output gate determines what information to output from the cell
state, and the hidden state (\(h_t\)) is computed as:
\[
o_t = \sigma(W_o \cdot [h_{t-1}, x_t] + b_o),
\]
\[
h_t = o_t \odot \tanh(C_t),
\]
where:
\begin{itemize}
    \item \(W_o\) is the weight matrix for the output gate,
    \item \(b_o\) is the bias term.
\end{itemize}

\section{The Role of the Cell State in Mitigating Vanishing Gradients}
The \textbf{cell state} (\(C_t\)) is the key mechanism that allows
LSTMs to mitigate the vanishing gradient problem. Here's why:

\subsection{Additive Updates to the Cell State}
The cell state is updated using \textbf{additive interactions} rather
than multiplicative ones:
\[
C_t = f_t \odot C_{t-1} + i_t \odot \tilde{C}_t.
\]
This additive nature ensures that gradients can flow through the cell
state without being multiplied by small values repeatedly, which
would cause them to vanish.

\subsection{Gradient Flow Through the Cell State}
During backpropagation, the gradient of the loss with respect to the
cell state (\(C_t\)) can flow through time with minimal attenuation.
This is because:
\[
\frac{\partial C_{t+1}}{\partial C_t} \approx \operatorname{diag}(f_{t+1}),
\]
(this is the direct, ``constant error carousel'' path; the gates also depend on \( h_t \), hence
indirectly on \( C_t \), which adds further terms that are usually neglected in this argument)
where \(f_{t+1}\) is the forget gate. Since \(f_{t+1}\) is typically
close to 1 (to retain information), the gradient does not vanish.

\subsection{Comparison with Standard RNNs}
In standard RNNs, the hidden state is updated multiplicatively:
\[
h_t = \sigma(W_h h_{t-1} + W_x x_t),
\]
which leads to vanishing gradients. In contrast, LSTMs use additive
updates to the cell state, which preserves gradient flow and allows
the network to learn long-term dependencies.

\section{Conclusion}
The \textbf{cell state} (\(C_t\)) in LSTMs is the primary reason why
LSTMs are effective at mitigating the vanishing gradient problem. By
using additive updates to the cell state, LSTMs allow gradients to
flow through time with minimal attenuation, enabling the network to
learn long-term dependencies effectively.
\end{notechapter}
